% GENERATED FILE. Edit canonical lessons, then run npm run generate:books.
% canonical-source-hash: fnv1a64:70665aa48010a895
% canonical-lessons: JA-C125-sugoi, JA-C125-hidoi, JA-C125-wakai, JA-C125-osanai, JA-C125-yutaka

\chapter{Amazing, Terrible, Young, Very young, Plentiful}
\label{ch:ja-amazing-terrible-young-very-young-plentiful}
\hlchaptermodality{125}

\begin{chapteropening}
\textbf{By the end of this chapter:} \emph{I can say amazing, terrible, young, very young and plentiful.}

Complete the last lesson of chapter 125: i can say amazing, terrible, young, very young and plentiful.
\end{chapteropening}

\section[\texorpdfstring{\ja{すごい} (sugoi)}{すごい (sugoi)}]{\ja{すごい} (sugoi) — amazing}

\label{lesson:JA-C125-sugoi}



\begin{quote}
Before the new one: say the Japanese for funny; strange, then the Japanese for interesting.
\end{quote}

\subsection*{You'll want to know: \ja{すごい}}
\textbf{\ja{すごい}} — \emph{sugoi} — \textquotedblleft{}amazing\textquotedblright{}.

\begin{grammarlens}[title={What you've built}]
One more word for this chapter.
\end{grammarlens}

\subsection*{Your turn}
\begin{itemize}
\raggedright
  \item \emph{Say it:} \emph{sugoi}
  \item \emph{Say it:} \emph{sugoi}, once more
  \item \emph{Say it:} say \emph{sugoi}
  \item \emph{Recall:} say the Japanese for funny; strange, then the Japanese for interesting, then say \emph{sugoi} again
\end{itemize}

\begin{tcolorbox}[breakable,colback=teal!4,colframe=teal!35!black,title={Before you move on}]
What does \ja{すごい} mean? (\textquotedblleft{}Amazing\textquotedblright{}.) Say it once more.
\end{tcolorbox}
\section[\texorpdfstring{\ja{ひどい} (hidoi)}{ひどい (hidoi)}]{\ja{ひどい} (hidoi) — terrible}

\label{lesson:JA-C125-hidoi}



\begin{quote}
Before the new one: say the Japanese for interesting, then the Japanese for amazing.
\end{quote}

\subsection*{You'll want to know: \ja{ひどい}}
\textbf{\ja{ひどい}} — \emph{hidoi} — \textquotedblleft{}terrible\textquotedblright{}.

\begin{grammarlens}[title={What you've built}]
One more word for this chapter.
\end{grammarlens}

\subsection*{Your turn}
\begin{itemize}
\raggedright
  \item \emph{Say it:} \emph{hidoi}
  \item \emph{Say it:} \emph{hidoi}, once more
  \item \emph{Say it:} say \emph{hidoi}
  \item \emph{Recall:} say the Japanese for interesting, then the Japanese for amazing, then say \emph{hidoi} again
\end{itemize}

\begin{tcolorbox}[breakable,colback=teal!4,colframe=teal!35!black,title={Before you move on}]
What does \ja{ひどい} mean? (\textquotedblleft{}Terrible\textquotedblright{}.) Say it once more.
\end{tcolorbox}
\section[\texorpdfstring{\ja{わかい} (wakai)}{わかい (wakai)}]{\ja{わかい} (wakai) — young}

\label{lesson:JA-C125-wakai}



\begin{quote}
Before the new one: say the Japanese for amazing, then the Japanese for terrible.
\end{quote}

\subsection*{You'll want to know: \ja{わかい}}
\textbf{\ja{わかい}} — \emph{wakai} — \textquotedblleft{}young\textquotedblright{}.

\begin{grammarlens}[title={What you've built}]
One more word for this chapter.
\end{grammarlens}

\subsection*{Your turn}
\begin{itemize}
\raggedright
  \item \emph{Say it:} \emph{wakai}
  \item \emph{Say it:} \emph{wakai}, once more
  \item \emph{Say it:} say \emph{wakai}
  \item \emph{Recall:} say the Japanese for amazing, then the Japanese for terrible, then say \emph{wakai} again
\end{itemize}

\begin{tcolorbox}[breakable,colback=teal!4,colframe=teal!35!black,title={Before you move on}]
What does \ja{わかい} mean? (\textquotedblleft{}Young\textquotedblright{}.) Say it once more.
\end{tcolorbox}
\section[\texorpdfstring{\ja{おさない} (osanai)}{おさない (osanai)}]{\ja{おさない} (osanai) — very young}

\label{lesson:JA-C125-osanai}



\begin{quote}
Before the new one: say the Japanese for terrible, then the Japanese for young.
\end{quote}

\subsection*{You'll want to know: \ja{おさない}}
\textbf{\ja{おさない}} — \emph{osanai} — \textquotedblleft{}very young\textquotedblright{}.

\begin{grammarlens}[title={What you've built}]
One more word for this chapter.
\end{grammarlens}

\subsection*{Your turn}
\begin{itemize}
\raggedright
  \item \emph{Say it:} \emph{osanai}
  \item \emph{Say it:} \emph{osanai}, once more
  \item \emph{Say it:} say \emph{osanai}
  \item \emph{Recall:} say the Japanese for terrible, then the Japanese for young, then say \emph{osanai} again
\end{itemize}

\begin{tcolorbox}[breakable,colback=teal!4,colframe=teal!35!black,title={Before you move on}]
What does \ja{おさない} mean? (\textquotedblleft{}Very young\textquotedblright{}.) Say it once more.
\end{tcolorbox}
\section[\texorpdfstring{\ja{ゆたか} (yutaka)}{ゆたか (yutaka)}]{\ja{ゆたか} (yutaka) — plentiful, rich}

\label{lesson:JA-C125-yutaka}



\begin{quote}
Before the new one: say the Japanese for young, then the Japanese for very young.
\end{quote}

\subsection*{You'll want to know: \ja{ゆたか}}
\textbf{\ja{ゆたか}} — \emph{yutaka} — \textquotedblleft{}plentiful, rich\textquotedblright{}.

\begin{grammarlens}[title={What you've built}]
One more word for this chapter.
\end{grammarlens}

\subsection*{Your turn}
\begin{itemize}
\raggedright
  \item \emph{Say it:} \emph{yutaka}
  \item \emph{Say it:} \emph{yutaka}, once more
  \item \emph{Say it:} say \emph{yutaka}
  \item \emph{Recall:} say the Japanese for young, then the Japanese for very young, then say \emph{yutaka} again
\end{itemize}

\begin{tcolorbox}[breakable,colback=teal!4,colframe=teal!35!black,title={Before you move on}]
What does \ja{ゆたか} mean? (\textquotedblleft{}Plentiful, rich\textquotedblright{}.) Say it once more.
\end{tcolorbox}
